\begin{figure*}[t]
\centering
\small
\setlength{\jot}{3pt}

\begin{tabular}{@{}p{0.36\textwidth}!{\color{black!20}\vrule width 0.4pt}p{0.58\textwidth}@{}}
\begin{minipage}[t]{\linewidth}
\vspace{0pt}
\centering
\textbf{Residual expression reduction}
\[
\begin{aligned}
  (\lambda x.\widehat e)\,\widehat v
  &\rightsquigarrow_{\mathrm{sym}}
  \widehat e[x\mapsto\widehat v],
  \\[2pt]
  \letin x{\widehat v}{\widehat e}
  &\rightsquigarrow_{\mathrm{sym}}
  \widehat e[x\mapsto\widehat v],
  \\[2pt]
  \avg D{\bar a}
  &\rightsquigarrow_{\mathrm{sym}}
  \operatorname{Mean}^{\#}_D(\bar a).
\end{aligned}
\]
\raggedright
The remaining reduction
rules are lifted from \cref{fig:semantics}. The affine expression
$\operatorname{Mean}^{\#}_D(\bar a)$ represents the primitive mean at the
symbolic arguments $\bar a$.
\end{minipage}
&
\begin{minipage}[t]{\linewidth}
\vspace{0pt}
\centering
\textbf{Symbolic small-step semantics}
\begin{mathpar}
  \inferrule*[right=Reduce]
    {\widehat e_1\rightsquigarrow_{\mathrm{sym}}\widehat e_2}
    {
      (\sigma\mid\widehat C[\widehat e_1]:\rho)
      \xrightarrow{}_{\mathrm{sym}}
      (\sigma\mid\widehat C[\widehat e_2]:\rho)
    }
  \and
  \inferrule*[right=Sample-E]
    {u\notin\operatorname{dom}(\sigma)}
    {
      (\sigma\mid\widehat C[\draw D\E{\bar a}]:\rho)
      \xrightarrow{}_{\mathrm{sym}}
      (\sigma,u\sim D(\bar a)\mid\widehat C[u]:\rho)
    }
  \and
  \inferrule*[right=Sample-G]
    {\operatorname{Domain}_D(\bar r)}
    {
      (\sigma\mid\widehat C[\draw D\G{\bar r}]:\rho)
      \xrightarrow{}_{\mathrm{sym}}
      (\sigma\mid\widehat C[r]:\rho)
    }
\end{mathpar}
\raggedright
The \textsc{Reduce} rule lifts reduction of residual symbolic expressions
to reduction of symbolic states under a symbolic evaluation context.
In \textsc{Sample-G}, the label $r$ is distributed according to
$D(\bar r)$. The label $\epsilon$ denotes a step that contributes no entry
to the $\G$ trace.
\end{minipage}
\end{tabular}

\caption{Symbolic semantics. Symbolic expression reduction
$\rightsquigarrow_{\mathrm{sym}}$ on residual expressions $\widehat e$ is
shown on the left, and the corresponding small-step transition relation
$\xrightarrow{}_{\mathrm{sym}}$ on symbolic states is shown on the right.
The \textsc{Reduce} rule lifts residual expression reduction through a
symbolic evaluation context, while \textsc{Sample-E} and \textsc{Sample-G}
handle the two sampling modes directly. Symbolic evaluation contexts
$\widehat C$ have the same grammar as the evaluation contexts $C$ in
\cref{fig:semantics}.}
\label{fig:symbolic-semantics}
\end{figure*}